% Replacement introduction for rvm_zeta_pi_mertens_v3.txt
% Insert after \tableofcontents and before the current section
% ``Purpose, separation of objects, and status''.
%
% Additional bibliography keys used here and not yet present in v3:
% Riemann1859, vonMangoldt1895, Beurling1937, BatemanDiamond1969,
% Diamond1977, DiamondMontgomeryVorhauer2006, BrouckeDebruyneVindas2020,
% Perelli2016, BenedettoPowellYilmaz2006, GunturkThao1995,
% DaubechiesDeVore2003, OdlyzkoTeRiele1985.

\section{Introduction}
\label{sec:introduction}

The classical theory of the Riemann zeta function joins four structures that
are usually introduced in different languages.  On the spectral side one has
the nontrivial zeros and their Riemann--von Mangoldt counting law; on the
analytic side one has the completed zeta function, its functional equation and
Hadamard product; on the prime side one has the prime-power measure and the
explicit formula; and on the inverse side one has the reciprocal Dirichlet
series and the Mertens summatory function.  Riemann's memoir made the first
three structures inseparable, and the later work of von Mangoldt, Hadamard and
Hamburger supplied, respectively, the rigorous zero-counting and explicit
formula, the entire-product factorization, and a converse characterization of
$\zeta$ from its Dirichlet-series and functional-equation data
\cite{Riemann1859,vonMangoldt1895,Hadamard1896,Hamburger1921,Titchmarsh1986,Edwards1974}.
The present paper asks what remains of that architecture when the fluctuating
zero staircase is replaced, at the outset, by the exact midpoint inverse of
its smooth Riemann--von Mangoldt clock.

The resulting construction is called the midpoint RVM system.  Its point of
departure is the smooth counting function
\begin{equation}
 \mathcal N_{\RVM}(T)
 :=\frac{T}{2\pi}\log\frac{T}{2\pi e}+\frac78,
 \label{eq:intro-clock}
\end{equation}
and the real ordinates $\gamma_k^{\RVM}$ determined by
\begin{equation}
 \mathcal N_{\RVM}(\gamma_k^{\RVM})=k-\frac12.
 \label{eq:intro-midpoint}
\end{equation}
The half-integer convention is not cosmetic.  It is the midpoint convention
for a staircase and leads to the exact Lambert inversion
\begin{equation}
 \gamma_k^{\RVM}
 =\frac{2\pi(k-11/8)}{W_0((k-11/8)/e)}.
 \label{eq:intro-lambert}
\end{equation}
Lambert-$W$ approximations to individual zero ordinates and smooth inversions
of the Riemann--von Mangoldt law occur in the literature, including explicit
equations and spectral approximations to the zeta zeros
\cite{DLMF-W,FrancaLeClair2014}.  What is used here is more specific: the
midpoint inverse in \eqref{eq:intro-midpoint} is treated as the exact spectrum
of a new deterministic system and is propagated, without inserting the
fluctuating term $S(T)$, through canonical products, explicit prime counting,
Mertens inversion and a continuous transport to the genuine ordinates.

The first RVM-space object is the canonical product
\begin{equation}
 P_{\RVM}^{\crit}(s)
 :=\prod_{k\ge1}
 \left(1+\frac{s(s-1)}{1/4+(\gamma_k^{\RVM})^2}\right),
 \qquad
 \zeta_{\RVM}(s):=\frac{P_{\RVM}^{\crit}(s)}{G(s)},
 \label{eq:intro-zeta-rvm}
\end{equation}
where
\begin{equation}
 G(s):=s(s-1)\Gamma(s/2)\pi^{-s/2}.
\end{equation}
Because $\gamma_k^{\RVM}\asymp k/\log k$, the paired factors form a normally
convergent genus-zero product in the symmetric variable $s(s-1)$.  Thus
$P_{\RVM}^{\crit}$ is entire and invariant under $s\mapsto1-s$, while
$\zeta_{\RVM}$ has a residue-one pole at $1$, the usual trivial zeros and the
prescribed simple nontrivial zeros
\begin{equation}
 \rho_k^{\RVM}=\frac12+i\gamma_k^{\RVM}.
\end{equation}
The location $\Re\rho_k^{\RVM}=1/2$ is an input of the synthetic model.  It is
not evidence that the genuine zero $\rho_k=\beta_k+i\gamma_k$ has
$\beta_k=1/2$.  Maintaining this distinction is essential: the paper constructs
an RVM analogue of the zeta--prime--Mertens triangle, not a relabelling of the
classical triangle.

There is a substantial surrounding literature on generalized prime systems.
In Beurling's theory the generalized primes may be arbitrary real numbers
$1<p_1\le p_2\le\cdots\to\infty$; their multiplicative semigroup defines
generalized integers and a zeta function by a Dirichlet series, an Euler
product or a Mellin--Stieltjes transform.  That framework shows both the power
and the flexibility of separating prime distribution from the ordinary
additive structure of the integers
\cite{Beurling1937,BatemanDiamond1969,Diamond1977}.  Later work exhibits
Beurling systems with prescribed analytic behaviour, critical zeros and large
oscillation, demonstrating that zeta-like zero data do not by themselves force
the classical arithmetic realization
\cite{DiamondMontgomeryVorhauer2006,BrouckeDebruyneVindas2020}.  The RVM system
is related to this tradition but is not introduced as an arbitrary Beurling
Euler product.  It starts from a prescribed zero clock, constructs an entire
completed product first, obtains a continuous prime-power count from its
explicit formula, and only afterwards asks whether that continuous object has
an atomic or integer-supported realization.

This reversal of direction is central.  From the RVM zeros we define the
prime-power counting function
\begin{equation}
 J_{\RVM}(x)
 :=\li(x)
 -2\Re\sum_{k\ge1}
 \Ei\!\left(\left(\frac12+i\gamma_k^{\RVM}\right)\log x\right)
 -\log2
 +\int_x^\infty\frac{dt}{t(t^2-1)\log t},
 \qquad x>1,
 \label{eq:intro-J}
\end{equation}
with the zero sum taken in midpoint-clock order.  The classical relation
between Riemann's prime-power count and the ordinary prime-counting function is
finite M\"obius inversion in the exponent variable
\cite{Riemann1859,Titchmarsh1986,Edwards1974}.  Applying precisely that
operation to \eqref{eq:intro-J} gives
\begin{equation}
 \pi_{\RVM}(x)
 :=\sum_{1\le m\le\log_2x}
 \frac{\mu(m)}{m}J_{\RVM}^{*}(x^{1/m}),
 \qquad x\ge2,
 \label{eq:intro-pi}
\end{equation}
and the exact inverse identity
\begin{equation}
 J_{\RVM}^{*}(x)
 =\sum_{1\le m\le\log_2x}\frac1m
 \pi_{\RVM}(x^{1/m}).
 \label{eq:intro-pi-inverse}
\end{equation}
The novelty intended here is not M\"obius inversion itself, which is classical,
but its use as the definition of the prime-counting coordinate in a spectral
system whose zeros were fixed before any prime atoms were supplied.

The function $\pi_{\RVM}$ is initially a continuous real-valued counting
coordinate, not a staircase.  It therefore supports two related notions of
``RVM prime''.  The first is measure-theoretic: the Stieltjes object
$d\pi_{\RVM}$ is the continuous prime measure generated by the RVM explicit
formula.  The second is a real level-crossing spectrum.  On every region where
$\pi_{\RVM}$ is continuous and strictly increasing we may define
\begin{equation}
 p_j^{\RVM}:=\pi_{\RVM}^{-1}(j),
 \label{eq:intro-real-primes}
\end{equation}
so that $\pi_{\RVM}(p_j^{\RVM})=j$.  These $p_j^{\RVM}$ are generally real and
need not be integers.  If global monotonicity is not available, the invariant
object is $d\pi_{\RVM}$ and the level-crossing notation is used only on
monotone branches.  This distinction prevents a continuous signed fluctuation
from being silently converted into a positive atomic prime system.

Sampling the real prime coordinate on unit cells produces
\begin{equation}
 c_n^{\RVM}
 :=\pi_{\RVM}(n)-\pi_{\RVM}(n-1),
 \qquad n\ge2.
 \label{eq:intro-cellweights}
\end{equation}
The number $c_n^{\RVM}$ is the RVM prime mass assigned to the cell $(n-1,n]$.
It is not assumed to be $0$ or $1$, and it need not initially be positive.
Consequently
\begin{equation}
 \zeta_{\mathbb Z,\RVM}(s)
 :=\prod_{n\ge2}(1-n^{-s})^{-c_n^{\RVM}}
 \label{eq:intro-weighted-euler}
\end{equation}
is a weighted analytic Euler object rather than the Euler product of an
ordinary set of integer primes.  This is the first continuous-to-atomic problem
of the paper: determine conditions under which the continuous real-prime
measure, its unit-cell masses, and an integer-supported weighted Euler product
represent the same analytic divisor up to a nonvanishing holomorphic factor.
The unit-periodic first harmonic predicted by the midpoint clock has exact
zero integral on every integer cell; the remaining question is therefore the
uniform size of the cell remainder.  Under the stated cell hypothesis one
obtains $c_n^{\RVM}\sim1/\log n$, but a positive $0/1$ realization is deliberately
not asserted without a separate selection mechanism.

The Mertens coordinate of RVM space is introduced independently of that
integer realization.  Analytically it is the symmetric inverse Mellin object
associated with
\begin{equation}
 \widehat M_{\RVM}^{\crit}(s)
 :=\frac1{\zeta_{\RVM}(s)}.
 \label{eq:intro-Mhat}
\end{equation}
A transition evaluation of the clock product, based on a smooth/sawtooth
Stieltjes decomposition and Abel regularization of the sawtooth Fourier series,
gives on fixed bands to the right of $1/2$
\begin{equation}
 \zeta_{\RVM}(s)
 =Z_{\RVM}(1+w(s))(1+O(1/|t|)),
 \qquad
 w(s)=\exp\{-2\pi i\mathcal N_{\RVM}(-i(s-1/2))\},
 \label{eq:intro-transition}
\end{equation}
and, on closed substrips of $1/2<\Re s<1$, the comparison
\begin{equation}
 \zeta_{\RVM}(s)
 =Z_{\RVM}(1+\chi(s))(1+O(1/|t|)).
 \label{eq:intro-transition-chi}
\end{equation}
It follows that $1/\zeta_{\RVM}(\sigma+it)$ tends to
$L=1/Z_{\RVM}$ for every fixed $\sigma>1/2$.  This vertical constant limit has
a consequence that sharply separates the RVM object from a classical Euler
system: neither $\zeta_{\RVM}$ nor its reciprocal is an absolutely convergent
ordinary Dirichlet series in any right half-plane.  Nevertheless, after
subtracting the constant vertical limit, the symmetric Perron inverse is
absolutely convergent and satisfies
\begin{equation}
 M_{\RVM}^{\crit}(x)
 =L+O_\varepsilon(x^{1/2+\varepsilon}).
 \label{eq:intro-Mbound}
\end{equation}
Thus the RVM Mertens function exists naturally as a continuous spectral inverse
even though no ordinary coefficient sequence for $1/\zeta_{\RVM}$ exists.
This point rules out a naive Dirichlet-convolution interpretation of the RVM
reciprocal and motivates a separate discretization.

The numerical and combinatorial companion to $M_{\RVM}^{\crit}$ is an
integer-valued walk obtained by first sampling a smooth target
$T(n)\approx M_{\RVM}^{\crit}(n)$.  In the implementation the leading
unit-periodic harmonic is removed before interpolation; it vanishes exactly at
integer arguments and therefore does not alter the prescribed integer samples.
The alphabet is restricted to $\{-1,0,+1\}$ on the squarefree support:
\begin{equation}
 \mu_h(n)=\mu(n)^2\varepsilon_n,
 \qquad \varepsilon_n\in\{-1,+1\},
 \qquad
 M_h(n):=\sum_{m\le n}\mu_h(m).
 \label{eq:intro-muh}
\end{equation}
The sign is assigned by the first-order feedback rule
\begin{equation}
 \varepsilon_n=
 \begin{cases}
  +1,&M_h(n-1)<T(n),\\
  -1,&M_h(n-1)\ge T(n),
 \end{cases}
 \qquad (\mu(n)^2=1),
 \label{eq:intro-quantizer}
\end{equation}
with $\mu_h(n)=0$ off the squarefree support.  This is a greedy
sigma--delta-type quantization: the next admissible increment is chosen from
$\{-1,+1\}$ according to the current tracking error.  Feedback quantization,
stable state variables and noise shaping are classical themes in sigma--delta
theory \cite{GunturkThao1995,DaubechiesDeVore2003,BenedettoPowellYilmaz2006};
the specific use made here is different.  The target is a zero-generated
Mertens inverse, the admissible times are selected arithmetically by
squarefreeness, and the output is studied as a Catalan excursion walk.  The
construction does not claim that $\mu_h$ is multiplicative or equal to the
classical M\"obius function.  Its purpose is to provide a controlled atomic
representative of the continuous RVM Mertens coordinate and to expose which
features survive discretization: tracking error, sign densities, correlation,
energy, spectral second moments and excursion geometry.

The squarefree support is natural but also restrictive.  Classically,
$\mu(n)=0$ exactly when $n$ is divisible by the square of a prime, while its
nonzero signs record the parity of the number of distinct prime factors.
Results on squarefree numbers in short intervals and on gaps between them show
that the support has rich and nontrivial local geometry; in the present paper
that geometry controls the admissible update times of the quantizer rather
than the location of the continuous real RVM primes.  The two discretizations
must therefore not be conflated: the cell masses $c_n^{\RVM}$ come from
$\pi_{\RVM}$, whereas the signs $\mu_h(n)$ are chosen dynamically to track
$M_{\RVM}^{\crit}$ on a prescribed support.

A further theme is comparison with the genuine arithmetic system.  At the
transform level the defect is exact:
\begin{equation}
 \mathcal Q_{\RVM}(s)
 :=\frac{\zeta(s)}{\zeta_{\RVM}(s)},
 \qquad
 \frac1{\zeta_{\RVM}(s)}
 =\mathcal Q_{\RVM}(s)\frac1{\zeta(s)}.
 \label{eq:intro-defect}
\end{equation}
No boundedness, invertibility or coefficient expansion of
$\mathcal Q_{\RVM}$ is built into this identity.  In particular,
\eqref{eq:intro-Mbound} is a theorem about the synthetic inverse and does not,
by itself, imply
$M(x)=O_\varepsilon(x^{1/2+\varepsilon})$.  The latter is a classical
equivalent formulation of the Riemann hypothesis, and the disproof of the
strong Mertens conjecture shows why the exponent and quantifiers must be stated
with care \cite{Titchmarsh1986,OdlyzkoTeRiele1985}.

To locate the exact endpoint obstruction, the midpoint ordinates are joined to
the genuine ordinates by
\begin{equation}
 \delta_k:=\gamma_k-\gamma_k^{\RVM},
 \qquad
 \gamma_{k,\kappa}:=\gamma_k^{\RVM}+\kappa\delta_k,
 \qquad 0\le\kappa\le1.
 \label{eq:intro-transport}
\end{equation}
The Riemann--von Mangoldt formula gives $\delta_k=O(1)$, so the associated
products
\begin{equation}
 P_\kappa(s)
 :=\prod_{k\ge1}
 \left(1+\frac{s(s-1)}{1/4+\gamma_{k,\kappa}^2}\right),
 \qquad
 \zeta_\kappa(s):=\frac{P_\kappa(s)}{G(s)}
 \label{eq:intro-path}
\end{equation}
converge jointly and normally.  Their moving divisor remains on the critical
line, and their logarithmic generator is an absolutely summable mode-by-mode
velocity.  The endpoint
\begin{equation}
 \zeta_{Z1}(s):=\zeta_1(s)
\end{equation}
has the genuine ordinates $\gamma_k$ but retains the imposed real part $1/2$.
It therefore exists unconditionally and must not be identified with $\zeta$
merely because its ordinates agree with those of the genuine zeros.

The analytic class of $\zeta_{Z1}$ makes the missing identification precise.
Its completed function has the zeta functional equation, residue one at
$s=1$, the usual trivial zeros and finite order.  Hamburger's converse theorem
belongs to the classical family of converse theorems that recover arithmetic
Dirichlet-series objects from analytic continuation, functional equations and
growth conditions \cite{Hamburger1921,Perelli2016}.  Applied here, it shows
that the remaining endpoint condition
\begin{equation}
 \tag{D}
 \zeta_{Z1}(s)=\sum_{n\ge1}\frac{a_n}{n^s}
 \quad\text{absolutely for }\Re s>1
 \label{eq:intro-D}
\end{equation}
is equivalent to
\begin{equation}
 \zeta_{Z1}\equiv\zeta
 \quad\Longleftrightarrow\quad
 \mathrm{RH}.
 \label{eq:intro-endpoint}
\end{equation}
Thus the transport does not prove RH; it isolates the continuous-to-atomic
closure that would do so.  In RVM language, the open endpoint question is
whether the zero-generated $Z1$ analytic object enters the ordinary discrete
Dirichlet category.  If it does, Hamburger uniqueness forces the classical
zeta function, and the associated prime and Mertens coordinates become the
classical ones.

The relationship among the principal RVM objects can therefore be summarized
as
\begin{equation}
 \{\gamma_k^{\RVM}\}
 \longrightarrow
 \zeta_{\RVM}
 \longrightarrow
 J_{\RVM}
 \longleftrightarrow
 \pi_{\RVM}
 \longrightarrow
 \bigl(d\pi_{\RVM},\,p_j^{\RVM},\,c_n^{\RVM}\bigr),
 \label{eq:intro-prime-chain}
\end{equation}
\begin{equation}
 \zeta_{\RVM}^{-1}
 \longrightarrow
 M_{\RVM}^{\crit}
 \longrightarrow
 (T,\mu_h,M_h),
 \label{eq:intro-M-chain}
\end{equation}
and
\begin{equation}
 \zeta_{\RVM}
 \longrightarrow
 \zeta_\kappa
 \longrightarrow
 \zeta_{Z1},
 \qquad
 \zeta_{Z1}=\zeta
 \Longleftrightarrow \mathrm{RH}.
 \label{eq:intro-transport-chain}
\end{equation}
The first chain is the prime coordinate, the second is the Mertens coordinate,
and the third is the ordinate-transport coordinate.  They share the same
midpoint spectrum but solve different realization problems.

\paragraph{Originality and scope.}
The ingredients surrounding the construction are classical or belong to
established neighbouring theories: the Riemann--von Mangoldt law, Lambert
inversion, Hadamard products, explicit formulas, M\"obius inversion, Perron
inversion, Beurling generalized primes, sigma--delta feedback and Hamburger
converse theorems.  The contribution claimed by this paper is their assembly
into a single midpoint RVM space with four coordinated levels:
\begin{enumerate}[label=\textup{(\roman*)}]
 \item an exact midpoint zero clock and its canonical zeta-like product;
 \item a continuous prime-power and prime coordinate, including real RVM
       prime levels and nonintegral unit-cell masses;
 \item a continuous synthetic Mertens inverse together with a squarefree,
       ternary, feedback discretization and its Catalan excursion geometry;
 \item a normal product transport from the midpoint spectrum to the genuine
       ordinates, ending in a single Dirichlet-series closure condition
       equivalent to RH.
\end{enumerate}
To the author's knowledge, the literature cited above contains the individual
background mechanisms but not this complete chain
\[
 \text{midpoint clock}
 \to \zeta_{\RVM}
 \to (J_{\RVM},\pi_{\RVM},p_j^{\RVM})
 \to (M_{\RVM}^{\crit},M_h)
 \to \zeta_{Z1}.
\]
This statement of originality is deliberately structural.  It does not claim
that every lemma or transform used in the chain is new, and it does not treat
numerical tracking as a proof of the endpoint condition.

\paragraph{Status of the results.}
The paper separates four logical levels.  The midpoint clock, canonical
product, exact divisor, logarithmic exponential-sum estimate, compact-uniform
explicit formula, finite M\"obius inversion, transition formula, vertical
limit, symmetric synthetic Mertens inverse, product transport and endpoint
analytic class are presented as theorems.  Square-root estimates for the RVM
prime coordinate and the local unit-cell remainder are stated with their
uniformity obligations visible.  The real-prime levels are defined on monotone
branches of $\pi_{\RVM}$, whereas a global positive real-prime sequence needs
global monotonicity.  The weighted Euler realization is conditional on
Mellin--Stieltjes and cell control.  The ternary walk $M_h$ is a discretization
and comparison object, not an Euler product and not the genuine Mertens
function.  Finally, condition \eqref{eq:intro-D} is the unresolved endpoint
criterion; proving it would identify the transported endpoint with the
classical zeta function and hence prove RH.

\paragraph{Organization.}
The paper begins with the midpoint Riemann--von Mangoldt clock, its Lambert
inversion and its bounded displacement from the genuine ordinates.  It then
constructs $\zeta_{\RVM}$ and proves the canonical-product and exponential-sum
results.  The following sections develop $J_{\RVM}$, $\pi_{\RVM}$, the real and
unit-cell prime coordinates, and the conditional weighted Euler realization.
The RVM Mertens transform and its transition asymptotics are treated next,
followed by the continuous-to-ternary discretization and its numerical and
excursion diagnostics where included.  The final analytic part transports the
midpoint ordinates to the genuine ordinates, proves the endpoint analytic
class, and formulates Hamburger condition (D) as the exact remaining closure.
